\section{Introduction}
\label{sec:introduction}

\subsection{Contexte}

\paragraph{L'agriculture de précision.} L'agriculture doit aujourd'hui produire davantage, pour une population qui continue de croître, avec moins d'eau, moins d'intrants et des sols souvent appauvris. L'agriculture de précision tente de répondre à cette équation en adaptant les interventions (irrigation, fertilisation, traitements) à la variabilité réelle de chaque parcelle, voire de chaque plant, au lieu d'appliquer la même dose partout \cite{zhang2002precision}. Gebbers et Adamchuk y voient l'un des leviers de la sécurité alimentaire, précisément parce qu'elle permet d'augmenter les rendements tout en réduisant les pertes et l'impact environnemental \cite{gebbers2010precision}. Cette approche repose sur des données (capteurs, images, positionnement par satellite) et sur des outils capables de les transformer en décisions ; c'est là que l'apprentissage automatique, puis l'apprentissage profond, ont trouvé leur place, de la prévision des rendements à la détection des adventices et des maladies \cite{liakos2018ml, kamilaris2018dl}.

\paragraph{Le poids économique des maladies des plantes.} Les maladies et les ravageurs comptent parmi les premières causes de pertes de récolte. À partir d'une évaluation mondiale fondée sur l'avis d'experts, Savary et al. estiment que les agents pathogènes et les ravageurs font perdre en moyenne 21,5\,\% de la récolte de blé, 30,0\,\% de celle de riz, 22,5\,\% pour le maïs, 17,2\,\% pour la pomme de terre et 21,4\,\% pour le soja, les pertes les plus lourdes touchant les régions déjà en déficit alimentaire \cite{savary2019burden}. Oerke arrivait à des ordres de grandeur comparables pour la période 2001--2003 : malgré les protections phytosanitaires en place, adventices, ravageurs et maladies réunis enlevaient de 26 à 29\,\% des récoltes de soja, de blé et de coton, et respectivement 31, 37 et 40\,\% de celles de maïs, de riz et de pomme de terre \cite{oerke2006crop}. Ces pertes ne se limitent pas aux rendements : elles fragilisent la sécurité alimentaire et pèsent d'abord sur les petits exploitants, qui ont le moins de moyens de les compenser \cite{strange2005plant}.

\paragraph{Les maladies des feuilles de canne à sucre.} La canne à sucre est, en tonnage, la culture la plus produite au monde : environ 1,9 milliard de tonnes en 2024, dont 39\,\% au Brésil et 23\,\% en Inde \cite{fao2025production}. Elle fournit environ 80\,\% du sucre produit dans le monde \cite{garsmeur2018mosaic}, mais aussi de l'éthanol et de la bagasse utilisée comme combustible. Sa culture est exposée à de nombreuses maladies fongiques, virales et bactériennes \cite{rott2000guide, viswanathan2011sugarcane}, et les quatre maladies étudiées ici ont toutes un coût économique documenté. En Inde, l'épidémie de pourriture rouge qui a frappé la variété Co~0238 a causé plus d'un milliard de dollars de pertes en une seule saison \cite{viswanathan2021redrot}. La mosaïque réduit la germination, le tallage, le nombre de tiges usinables, leur poids et leur teneur en saccharose \cite{viswanathan2005mosaic}. La rouille brune a fait baisser de 16\,\% le tonnage de cannes et de 14\,\% le rendement en sucre de la variété LCP~85-384 \cite{hoy2009brownrust}, et la rouille orangée réduit nettement la photosynthèse des variétés sensibles \cite{zhao2011orangerust}. Quant à la maladie de la feuille jaune, due au virus ScYLV, elle a réduit de 28\,\% le poids des tiges et de 11\,\% la teneur en sucre sur la variété R577, et le rendement de 37\,\% en première repousse \cite{rassaby2003yellow, elsayed2015scylv}. Dans tous les cas, une détection précoce permet de limiter la propagation : arracher les touffes atteintes, éviter de prélever des boutures sur des parcelles contaminées ou orienter le choix des variétés.

\paragraph{Les limites du diagnostic traditionnel.} Le diagnostic repose encore le plus souvent sur l'inspection visuelle des parcelles par un agronome ou un technicien. Cette méthode est lente, difficile à étendre à de grandes surfaces, et dépend de l'expérience de l'observateur. Dans leur revue des méthodes d'estimation de la sévérité, Bock et al. relèvent une tendance générale à surestimer la maladie lorsqu'elle est encore peu étendue (moins de 10\,\% de la surface), précisément au stade où l'intervention serait la plus utile, et notent que la concordance entre évaluateurs, voire chez un même évaluateur d'une fois sur l'autre, peut varier sensiblement en l'absence de formation ou d'échelles de référence \cite{bock2010severity}. Les techniques de laboratoire (tests sérologiques, PCR) sont fiables, mais elles sont coûteuses, demandent des prélèvements et un délai d'analyse, et ne se prêtent pas à un suivi régulier au champ \cite{martinelli2015advanced}. L'imagerie numérique offre un compromis intéressant entre ces deux extrêmes, à condition de disposer d'algorithmes capables d'interpréter les images de manière fiable \cite{mahlein2016plant, barbedo2016review}.

\paragraph{L'intérêt de la vision par ordinateur et du deep learning.} Depuis la victoire d'AlexNet au concours ImageNet en 2012 \cite{krizhevsky2012imagenet}, les réseaux de neurones convolutifs (CNN) se sont imposés en vision par ordinateur \cite{lecun2015deep}. Au lieu de concevoir à la main des descripteurs de couleur ou de texture, le réseau apprend directement, à partir des pixels, une hiérarchie de caractéristiques adaptée à la tâche. Appliqués aux maladies des plantes, ils ont dépassé 99\,\% d'accuracy sur la base PlantVillage \cite{mohanty2016using, ferentinos2018deep}. Associés à des smartphones de plus en plus répandus dans les campagnes, ils rendent envisageable un diagnostic immédiat, au champ, par l'agriculteur lui-même.

\subsection{Problématique}

Ces bons résultats ont une contrepartie bien connue : un CNN moderne compte des dizaines de millions de paramètres et a besoin de beaucoup d'exemples annotés pour les ajuster. ImageNet-1k, sur lequel la plupart des architectures sont entraînées, compte 1\,281\,167 images d'apprentissage réparties en 1\,000 classes \cite{russakovsky2015imagenet}, et Sun et al. ont montré que les performances continuent de progresser, de façon logarithmique, avec le volume de données d'apprentissage \cite{sun2017revisiting}.

Les jeux de données agricoles sont, eux, presque toujours petits. Photographier des feuilles malades au bon stade, dans des conditions variées, puis les faire étiqueter par des spécialistes prend du temps et coûte cher. Le jeu utilisé dans ce projet ne compte que 2\,521 images, 2\,287 après nettoyage, soit environ cinq cents fois moins qu'ImageNet. Barbedo a d'ailleurs montré que la taille et surtout la diversité des données restent le principal frein à la généralisation des modèles de diagnostic \cite{barbedo2018impact}. Entraîner un réseau profond à partir de zéro dans ces conditions mène rapidement au surapprentissage.

L'apprentissage par transfert apporte une réponse : réutiliser un réseau entraîné sur un grand corpus généraliste et l'adapter au nouveau problème \cite{pan2010survey}. Les premières couches d'un CNN apprennent des motifs génériques (contours, textures, couleurs) qui se transfèrent bien d'un domaine à l'autre \cite{yosinski2014transferable}, et, sur 16 réseaux et 12 jeux de données, Kornblith et al. trouvent une corrélation très forte entre la précision sur ImageNet et la précision après transfert ($r=0{,}99$ lorsque le réseau sert d'extracteur fixe, $r=0{,}96$ après fine-tuning) \cite{kornblith2019better}. Plusieurs questions restent cependant ouvertes, et ce sont elles qui guident ce projet :
\begin{itemize}
  \item jusqu'où faut-il adapter le réseau pré-entraîné : geler toutes les couches et n'apprendre qu'un classifieur (extraction de caractéristiques), ré-entraîner seulement les couches profondes (fine-tuning partiel), ou tout ré-entraîner (fine-tuning complet) ?
  \item quelle architecture tire le meilleur parti de si peu de données, et à quel coût de calcul ?
  \item combiner les représentations de plusieurs réseaux pré-entraînés apporte-t-il un gain réel ?
  \item les décisions des modèles reposent-elles vraiment sur les symptômes visibles sur la feuille ?
\end{itemize}

\subsection{Contributions}

Ce travail apporte les contributions suivantes :
\begin{itemize}
  \item une \textbf{comparaison de quatre architectures pré-entraînées} récentes et de complexités différentes (EfficientNetV2-S, MobileNetV4-Conv-M, ConvNeXtV2-T et SE-ResNet-50), confrontées à un CNN entraîné from scratch, dans des conditions d'entraînement strictement identiques ;
  \item une \textbf{étude de l'impact du fine-tuning}, à travers trois stratégies de transfert (extraction de caractéristiques, fine-tuning partiel et fine-tuning complet) appliquées à chaque architecture ;
  \item une \textbf{étude de la fusion de caractéristiques}, évaluée de façon exhaustive sur les 15 combinaisons possibles des quatre réseaux, et comparée à une fusion des décisions par vote souple ;
  \item une \textbf{analyse des performances sur les différentes maladies}, complétée par une évaluation de la robustesse des modèles face à des dégradations d'image et de leur coût de calcul ;
  \item l'\textbf{interprétation des prédictions par Grad-CAM}, pour vérifier que les modèles regardent bien les zones malades ;
  \item un \textbf{protocole d'évaluation soigné} : audit du jeu de données, qui a révélé 210 doublons exacts et 24 quasi-doublons, ensemble de test isolé et validation croisée stratifiée à 5 plis.
\end{itemize}

\subsection{Organisation du document}

La suite du rapport est organisée comme suit. La section~\ref{sec:etat-art} présente l'état de l'art, des méthodes classiques de traitement d'images aux travaux récents sur la canne à sucre, et situe notre travail par rapport à eux. La section~\ref{sec:methodologie} détaille la méthodologie : jeu de données et nettoyage, prétraitement, augmentation, modèles, stratégies de transfert, fusion, paramètres d'entraînement et métriques. La section~\ref{sec:resultats} présente les résultats et compare les modèles selon leur précision, leur légèreté et leur robustesse. La section~\ref{sec:interpretabilite} analyse les cartes Grad-CAM. La section~\ref{sec:discussion} discute les forces et les limites de l'étude et propose des perspectives, avant la conclusion de la section~\ref{sec:conclusion}.
